\section{Query-Time Set Selection}
\label{sec:method}

Figure~\ref{fig:overview} summarizes the design. We first formulate retrieval as returning a set of rounds and define the two measures that the rest of the paper uses, full-set recall for a retrieved set and index visibility for an index (Section~\ref{sec:formulation}). The memory then writes a compact index at ingestion (Section~\ref{sec:index}), a selector reads this index at question time and returns the evidence set (Section~\ref{sec:select}), and a reader answers from the selected raw rounds (Section~\ref{sec:read}). The remaining subsections address the limitations that the evaluation reveals: selectors that are too weak (Section~\ref{sec:light}), indexes that hide the evidence and readers that fail to aggregate it (Section~\ref{sec:vgs}), and the cost of the strong selection call (Sections~\ref{sec:deploy} and~\ref{sec:cost}).

\subsection{Problem Formulation}
\label{sec:formulation}

A memory holds the history of one user as a sequence of rounds $H=(r_1,\dots,r_n)$. Each round $r_i$ contains a user message, an assistant reply, zero or more images, a date, and the identifier of its session. Given a question $q$, the memory returns a subset $S(q)\subseteq H$, and a reader answers $q$ from the rounds in $S(q)$. The benchmarks annotate a gold evidence set $G(q)\subseteq H$, the rounds that contain the decisive evidence. We measure retrieval by full-set recall,
\begin{equation}
\mathrm{FSR}(q)=\mathbb{1}\left[G(q)\subseteq S(q)\right],
\end{equation}
which equals one only when every gold round is returned. Averaged over questions, FSR is the per-question form of the all-recall measure used for multi-hop retrieval~\citep{gutierrez2024hipporag}, and it plays the role that supporting-fact exact match plays in multi-hop question answering~\citep{yang2018hotpotqa}. We report it together with the coverage $|G(q)\cap S(q)|/|G(q)|$ and the size $|S(q)|$. Because the reader pays for every opened round, a retriever should reach a high full-set recall with a small $|S(q)|$. Top-$k$ retrieval fixes $|S(q)|=k$ and ranks rounds by $\cos(e(q),e(r_i))$, so its full-set recall is determined by the rank of the worst-ranked gold round; on MemEye, this rank has a median of 5 and exceeds 40 for 7.8\% of questions.

Index-based selection can only find evidence that its index exposes. We therefore measure index visibility: for every gold round, the fraction of the content words that the round shares with the question and that also appear in its index line, averaged over gold rounds. The measure needs no model call and can be computed for any benchmark before a retrieval design is chosen.

\subsection{Memory Index}
\label{sec:index}

To make the whole history readable at question time, the memory writes one index line per round at ingestion. The line contains the round identifier and date, the first 200 characters of the user message, a one-sentence description of every image produced by a captioner, and the first 100 characters of the assistant reply. On MemEye, an index line costs about 70 tokens, so the index of a 300-round history fits in about 22k tokens, and a context of one million tokens holds the index of roughly 14,000 rounds. The index is written once, is never shown to the reader in place of the images, and changes only by appending, so its tokens can be served from a cache across questions.

\subsection{Selection at Question Time}
\label{sec:select}

Given the index of $H$ and the question $q$ (with its image, if any), a selector model returns a list of round identifiers. The instruction asks the selector to include every round that may hold evidence, in particular all members of any group that the question refers to (for example, all advertisements of a brand or every version of a screen), the rounds that define labels used by the question, and rounds whose image description is too coarse to rule them out (Appendix~\ref{app:prompts}). The output is capped at $N{=}40$ identifiers but is otherwise free, so the size of $S(q)$ is decided per question. This set formulation is the main difference from listwise reranking~\citep{sun2023rankgpt}, which orders candidates and leaves the cut-off to a fixed $k$, and from index-based reading agents~\citep{lee2024readagent,chen2023memwalker}, which look up passages of one long document.

To isolate the set formulation, we compare it with controls that use the same model, candidates, and index lines. The ranking control asks the model to rank all candidates by usefulness and keeps, for every question, as many top-ranked rounds as \qss{} selected, so that it receives the set size for free. Because the instruction of \qss{} also asks for every possible member, a second ranking control uses this recall-oriented instruction and changes only the output, from an unordered set to a complete ranking with the same cut. The dynamic cut-off control asks the model to grade every candidate from 0 (not relevant) to 3 (contains needed evidence) in one listwise call, in the manner of graded-label rankers~\citep{zhuang2024beyond}, and keeps the candidates graded 2 or higher; we fixed this threshold before any run.

For long histories, \qss{} first ranks rounds by cosine similarity and shows the selector only the index lines of the top-$P$ candidates in chronological order ($P{=}100$ for the original histories and $P{=}200$ for the enlarged ones). The prefilter bounds the selection cost by $P$ rather than $n$, while the selector still decides membership jointly over the candidates.

\subsection{Reading the Selected Set}
\label{sec:read}

The reader receives the raw rounds of $S(q)$ in chronological order, with images at its standard resolution. Isolated rounds, however, lack the narrative structure that questions such as ``the first update'' or ``the D9 summary'' presuppose, namely the session order, the labels, and the dates of the surrounding history. In the index-as-context variant, the reader therefore receives the full index before the selected rounds, so that it can locate them in the history while still answering from the original images.

\subsection{Lightweight Selectors}
\label{sec:light}

A lightweight selector reads the same index but, as Section~\ref{sec:rq4} shows, chooses incomplete or wrong sets. We study two remedies that keep every question-time call on the lightweight model. Sampling consensus draws $K{=}5$ selections at temperature 1 and keeps the rounds chosen by at least $V$ samples; we use $V{=}2$, which trades part of the recall gain of the union ($V{=}1$) for the precision that a lightweight reader needs. Resolution write-back stores, after a question has been answered, the pair of the question and its selected set $S(q)$. Later selections for the same user see the stored resolutions below the index, so that a reference resolved once, such as ``the D9 summary card'', can be reused. In our experiments, the stored sets come from the strong selector, questions arrive in a random order, each question sees only resolutions of questions that precede it, and resolutions of near-paraphrases (question-embedding cosine above 0.9) are withheld, so that the evidence set of a question cannot leak through a reworded copy. A model-free variant reuses the selected sets of the three most similar resolved questions together with the cosine top-5 rounds.

\subsection{Visibility-Gated Set Reading}
\label{sec:vgs}

The held-out evaluation of Section~\ref{sec:rq3} shows that index-based selection fails when the index hides the evidence, and the failure analysis of Section~\ref{sec:rq5} shows that lightweight readers fail to aggregate evidence sets even when the sets are complete. Visibility-gated set reading (VGS) addresses both failures with one decision per question. To make this decision without labels, VGS estimates index visibility at question time: for the ten rounds with the highest cosine similarity, it computes the share of the question words found in each round that its index line also shows, and it averages these shares into $\hat{v}(q)$. When $\hat{v}(q)\ge\tau$, the index exposes the evidence, and VGS runs \qss{} with the standard index and standard reading. When $\hat{v}(q)<\tau$, VGS widens the index lines to 600 characters of the user message for selection, and it adds per-member notes for reading: a lightweight model reads each selected round in full, including its images, and writes down the facts that bear on the question, and the reader receives these notes before the raw rounds. The notes make the members of the set and their values explicit, so that the reader aggregates over a short list instead of searching long messages; because the raw rounds remain in the input, a note that misses a detail does not remove it. We fixed $\tau{=}0.7$ from the development data, where the estimate separates short-message histories (98.7\% of MemEye questions above 0.7) from long-message histories (12.4\% of the MemLens development questions above 0.7).

\subsection{Cost-Aware Deployment}
\label{sec:deploy}

The strong selection call dominates the cost of \qss{} and VGS: it reads the whole index for every question, and its output alone, including the minimum thinking budget, costs about as much as reading 40 rounds with the lightweight reader. We study three deployment mechanisms that leave the decision of the selector unchanged. Set-aware routing classifies each question from its text alone, with one call of the lightweight model, as single-item or multi-item, that is, according to whether answering it requires combining several items, events, or points in time; single-item questions go to top-$k$ retrieval, and only multi-item questions pay for set selection and, for the lightweight reader, a thinking budget. A cache-native index exploits the fact that the index of a persistent memory changes only by appending: the system instruction and the index are stored once per user in an explicit context cache of the API, and every selection call, as well as the index-as-context reader, sends only the question and the opened rounds. Cached tokens are billed at a tenth of the input price, plus storage per hour. A compact index, finally, rewrites each line deterministically to reduce the tokens that every call reads: it drops the assistant snippet, lists rounds by short positional identifiers under one header per session (keeping the session labels of the benchmark, which questions use as references), and removes function words from the user text before truncation, so that the same character budget shows more of the message.

\subsection{Cost Accounting}
\label{sec:cost}

A selection call reads text only, namely the index or the prefiltered index and the question, and it writes a short list of identifiers. On MemEye, a call of the strong selector reads 11.5k tokens on average and writes 230 tokens including thinking, which costs \$9.47 per 1,000 questions at list prices without caching. With the explicit context cache of Section~\ref{sec:deploy}, the measured cost of this call falls to \$1.76, plus \$0.28 for creating and storing the caches. The reader pays only for the opened rounds, 7.3 on average for \qss{} instead of 40 for top-40 retrieval. We report measured costs throughout. For the development runs of Section~\ref{sec:rq2}, which did not use the explicit cache, we additionally give warm-cache estimates, and the measured cached runs of Section~\ref{sec:rq6} lie within 8\% of them; Appendix~\ref{app:cost} gives the cost model.
